\subsection{Auxiliary finite-model illustration}

The following exact construction isolates the distinction between operating
profiles and generative modules in the main paper's finite model.  It is an
illustration rather than an additional assumption or theorem.

\begin{example}[A two-module generator]\label{ex:modular-generator}
Let \(X=\{\mathrm{down},\mathrm{up}\}\), let
\(\mu_{b^-}=(3/4,1/4)\) and \(\mu_{b^+}=(1/4,3/4)\), and take identity
closure.  Strategies \(s_T\) and \(s_H\) carry modules \(m_T\) and \(m_H\).
A project \(q\) requires both.  When
\(\{m_T,m_H\}\subseteq C\), set
\[
  G(q,b,C)(\operatorname{some}(c))=\frac23,\qquad
  G(q,b,C)(\operatorname{none})=\frac13,\qquad
  \nu(q,b,C,c)=\frac34;
\]
otherwise set \(G(q,b,C)=\delta_{\operatorname{none}}\).  Retaining both
modules therefore admits \(c\) with probability \(1/2\), while deleting the
sole carrier of either module makes admission impossible.  If
\(u_c(\mathrm{down})=-1\) and \(u_c(\mathrm{up})=2\), then
\(j_c(b^-)=-1/4\) and \(j_c(b^+)=5/4\): current operating payoff depends on
the filtered state, while candidate availability depends on retained closure.
For a positive-duration instance, take \(d_q=2\), initiation cost \(1/10\),
and operation flag \(o_q=1\).
\end{example}

\subsection{Conditional replacement under a capacity bound}
\label{app:conditional-replacement}

The supporting question is what an incumbent should release when an
outer-certified candidate \(c\notin L\) becomes available for retention and
capacity binds.  This formulation is conditional on outer certification and
is separate from the raw generation, verification, and admission law in the
main paper.

Let the incumbent satisfy \(W(L)\le B\), let \(w_c\le B\), and permit deletion
of active incumbents \(D\subseteq L\setminus\{s_0\}\).  Additivity makes the
required release
\[
 r_c(B):=[W(L)+w_c-B]_+,
\]
so the feasible deletion family is
\[
 \mathcal D_c(B)
 :=\{D\subseteq L\setminus\{s_0\}:W(D)\ge r_c(B)\}.
\]
The conditional replacement problem is
\begin{equation}\label{eq:optional-replacement-problem}
 \operatorname*{maximize}_{D\in\mathcal D_c(B)}
 V_\theta\bigl(b,(L\setminus D)\cup\{c\}\bigr).
\end{equation}

For interpretation, define the unconstrained candidate gain and displacement
loss
\begin{align*}
 G_c(b,L)&:=V_\theta(b,L\cup\{c\})-V_\theta(b,L),\\
 \ell_c(D)&:=V_\theta(b,L\cup\{c\})
 -V_\theta\bigl(b,(L\setminus D)\cup\{c\}\bigr).
\end{align*}
Then the best net admission value is the definition-level identity
\[
 A_c(b,L,B)=G_c(b,L)-\min_{D\in\mathcal D_c(B)}\ell_c(D).
\]
Under common insertion and the productive factorization, a deletion safe
before admission is sufficient for zero displacement loss, but it is not
necessary: the candidate may replace a capability carried by a deleted
incumbent.  This supporting formulation is not part of the Lean-verified
theorem set, and it makes no constrained--penalized equivalence claim.
